\section{Approche réflexive}

\subsection{Si vous deviez ne retenir qu’une seule notion de ce travail, quelle serait-elle ?}
Si je ne devais retenir qu’une seule notion de ce travail, ce serait l’importance cruciale d’une approche holistique et centrée sur la personne lors de l’initiation d’un traitement. L’EBM ne se limite pas aux statistiques d’efficacité brute, elle doit intégrer les comorbidités et, surtout, le ressenti subjectif du patient. Considérer une plainte qui impacte le quotidien est le garant de l’alliance thérapeutique. Sans ce sentiment d’être écoutée, la compliance est menacée, particulièrement dans des pathologies instables comme la dépression mixte.

Cette vision rejoint la pensée d’Édouard Zarifian qui, dans La Force de guérir \parencite{zarifian1999force}, souligne que la pratique psychiatrique est indissociable du lien humain et du contexte culturel. Il nous rappelle que la prise en charge s'inscrit dans un cadre social large qui influence directement le regard des médecins et leurs processus diagnostiques. L'élimination du seul symptôme par la technique est insuffisante si l'on ignore le sens que le patient lui donne au sein de son histoire personnelle.

\subsection{Que pensez-vous avoir bien réussi dans ce travail ?}
Le rapport de cas a été un exercice fréquent lors de mon stage, à travers celui-ci j’ai tenté de m’améliorer au mieux. Concrètement, j'ai tenté de décrire un tableau global tout en insistant sur les éléments pertinents. La patiente a vécu beaucoup d’expériences parfois traumatisantes dans sa vie. Plusieurs interprétations peuvent être faites sur l’évolution de la patiente et sur les raisons de son hospitalisation. Grâce à la collaboration étroite avec les psychiatres présents dans le service, j’ai pu communiquer mes pensées et démarches cliniques afin qu’ils puissent me guider et me corriger si nécessaire. Il y également eu un apprentissage constant pour réaliser une transmission claire et des hypothèses construites. 

\subsection{Que pensez-vous devoir améliorer ?}
J’ai tenté de maîtriser la recherche sur différentes bases de données au mieux mais je ressens un manque d’expérience particulièrement pour le choix des MESH-Terms. L’interprétation de chaque donnée individuelle a été, selon moi, bien comprise et réalisée mais mon assurance diminue quand la décision finale de la fiabilité de l’article n’est pas évidente et reste nuancée. J'ai tenté de m’appuyer sur des éléments objectifs afin d’aiguiser mon opinion.  

\subsection{Quelle(s) question(s) souhaiteriez-vous approfondir dans un prochain travail ?}
Je pense qu’aborder le suivi plus long terme de la patiente et son évolution dans le temps aurait pu être très intéressant. En outre, la démarche sociale et psychologique aurait pu être davantage approfondie et mise en avant. Effectivement, je suis convaincue que la prise en charge doit être multidisciplinaire et que l’aspect purement médicamenteux de cette prise en charge est nécessaire mais loin d'être suffisante et la place de la psychothérapie a une place déterminante dans l’évolution de l’état psychique de la patiente. Je pense que se renseigner sur les différents types de psychothérapie -- leur durée, les effets de la communication sur la patiente -- est primordial à maîtriser. Sans doute, dans notre médecine actuelle, le focus est peut-être trop axé sur la pharmacologie. 